\documentclass[10pt]{article}
\usepackage[margin=1in]{geometry}
\usepackage[T1]{fontenc}
\usepackage[utf8]{inputenc}
\usepackage{amsmath,amssymb,amsthm}
\usepackage{enumitem}
\usepackage{hyperref}

\newtheorem*{theorem}{Theorem}
\newcommand{\PSh}{\mathsf{PSh}}
\newcommand{\PtSh}{\mathsf{PtSh}}
\newcommand{\SFin}{\mathsf{SFin}}
\newcommand{\Cart}{\mathsf{Cart}}
\newcommand{\Vcat}{\mathsf{V}}

\title{A Point-Generated Linear--Nonlinear Sheaf Bridge\\
\large Extended abstract, AI-assisted draft}
\author{Hoapa project}
\date{June 2026}

\begin{document}
\maketitle

\begin{abstract}
We construct, in Cubical Agda, a sheaf-level linear--nonlinear adjunction over a finite point-generated actor-causal site.  The site is restricted to supported finite frames: every actor appearing in a frame is incident to some event.  Its topology is generated by the canonical cover of a frame by its single-event point subframes.  For any value category with finite products, the resulting pointification functor sends a presheaf to the product of its event-local values; this functor is proved to be left adjoint to the inclusion of point sheaves into supported presheaves.  Given a value-level Benton linear--nonlinear adjunction and finite products on the linear value category, this reflector removes the sheafification hypothesis from the point-generated sheaf bridge and yields a machine-checked adjunction between cartesian-valued and linear-valued point sheaves.
\end{abstract}

\section*{The mathematical statement}

Let $A$ be an actor base, and let $\SFin_A$ be the category of supported finite actor-causal frames.  Objects are finite local frames whose actors are all supported by events; morphisms are the strict frame embeddings used by the actor-causal site, restricted to supported frames.  The point topology on $\SFin_A$ is generated by the canonical family
\[
  \{\operatorname{pt}_e(F) \longrightarrow F\}_{e \in E(F)},
\]
where $\operatorname{pt}_e(F)$ is the one-event subframe incident at $e$.  This is deliberately the event-discrete topology: it records the universal reconstruction of a frame from its event-local views, while leaving causal and higher-overlap refinements for later topologies.

For a category $\mathcal C$ with finite products, write
\[
  \PSh_A(\mathcal C) = [\SFin_A^{\mathrm{op}},\mathcal C]
\]
for the implemented category of supported presheaves, and let $\PtSh_A(\mathcal C) \subseteq \PSh_A(\mathcal C)$ be the full subcategory of point sheaves.  The point sheaf condition is the descent condition for the canonical point cover; equivalently, a presheaf is determined by its event-local values.  The formalization proves this equivalence as a comparison between point descent and a product-locality condition.

The reflector is explicit.  For $P \in \PSh_A(\mathcal C)$ define
\[
  a_{\mathcal C}(P)(F) = \prod_{e\in E(F)} P(\operatorname{pt}_e(F)).
\]
Restrictions are induced by the corresponding point-to-point morphisms.  The proof that this is functorial, point-local, and universal is carried out behind an abstract finite-product interface, so that Agda does not unfold products into enumeration-dependent normal forms during unification.  The machine-checked result is an adjunction
\[
  a_{\mathcal C} \dashv \iota_{\mathcal C}
  \qquad
  (\iota_{\mathcal C}:\PtSh_A(\mathcal C)\hookrightarrow\PSh_A(\mathcal C)).
\]

Now suppose that $\Cart$ is a cartesian value category with finite products, that $\Vcat$ is the category underlying a star-autonomous value structure and also has finite products, and that we are given a Benton-style value-level adjunction
\[
  F \dashv G : \Cart \rightleftarrows \Vcat .
\]
Postcomposition lifts $F$ and $G$ to supported presheaf categories:
\[
  F_* : \PSh_A(\Cart) \to \PSh_A(\Vcat),
  \qquad
  G_* : \PSh_A(\Vcat) \to \PSh_A(\Cart),
\]
with a pointwise adjunction $F_*\dashv G_*$.  The right adjoint preserves point sheaves: descent along the point cover is transported across the natural bijection
\[
  \Cart(X, GY) \cong \Vcat(FX,Y).
\]
The left adjoint is made sheaf-level by reflection.  Hence define
\[
  R_{\mathrm{pt}} = G_*\vert_{\PtSh_A(\Vcat)},
  \qquad
  L_{\mathrm{pt}} = a_{\Vcat}\,F_*\,\iota_{\Cart}.
\]

\begin{theorem}
Assume $\Cart$ has finite products, $\Vcat$ has finite products, and $F\dashv G:\Cart\rightleftarrows\Vcat$ is a value-level linear--nonlinear adjunction.  Then there is a sheaf-level adjunction
\[
  L_{\mathrm{pt}} \dashv R_{\mathrm{pt}} :
  \PtSh_A(\Cart) \rightleftarrows \PtSh_A(\Vcat)
\]
over the supported point-generated actor-causal site.  In particular, the linear exponential at this sheaf level is the comonad $L_{\mathrm{pt}}R_{\mathrm{pt}}$.
\end{theorem}

The result is formalized without postulates in the new modules \texttt{SupportedChangeOfBase.agda}, \texttt{SupportedLinearNonlinearSheafAdapter.agda}, and \texttt{PointLinearNonlinearAdjunction.agda}, building on the point reflector constructed in \texttt{PointSheafAdjunction.agda}.

\section*{Scope and significance}

The theorem is modest but precise.  It is a Grothendieck-style descent result over a finite site, a Lawvere-style use of adjointness as structure, and a Benton-style linear--nonlinear bridge lifted from values to sheaves.  The point-generated topology is intentionally small: it proves that event-local observations freely generate a coherent sheaf-level object and that this reconstruction is compatible with a linear--nonlinear value adjunction.

Equally important are the limitations.  The theorem does not construct the value-level adjunction $F\dashv G$; it assumes it.  It does not yet instantiate the linear value category with a concrete sessions model satisfying the finite-products hypothesis.  It is not a full sheafification theorem for all finite actor-causal covers, nor a comparison with a Cubical library coverage for the point topology.  Finally, it is not the cross-site adjunction to the modal side.  It is the first checked bridge on the operational side: point-generated sheafification plus a sheaf-level linear--nonlinear adjunction.

\begin{thebibliography}{9}

\bibitem{SGA4}
M. Artin, A. Grothendieck, and J.-L. Verdier, eds.
\emph{Th\'eorie des topos et cohomologie \'etale des sch\'emas (SGA 4)}.
Lecture Notes in Mathematics 269, 270, 305. Springer, 1972--1973.

\bibitem{Lawvere1969}
F. W. Lawvere.
Adjointness in foundations.
\emph{Dialectica} 23(3--4):281--296, 1969.

\bibitem{MacLaneMoerdijk1992}
S. Mac Lane and I. Moerdijk.
\emph{Sheaves in Geometry and Logic: A First Introduction to Topos Theory}.
Springer, 1992.

\bibitem{Benton1995}
P. N. Benton.
A mixed linear and non-linear logic: proofs, terms and models.
In \emph{Computer Science Logic}, Lecture Notes in Computer Science 933, pp. 121--135. Springer, 1995.

\bibitem{CattaniWinskel1996}
G. L. Cattani and G. Winskel.
Presheaf models for concurrency.
BRICS Report Series RS-96-35, 1996.

\bibitem{NielsenWinskel1995}
M. Nielsen and G. Winskel.
Models for concurrency.
In S. Abramsky, D. M. Gabbay, and T. S. E. Maibaum, eds., \emph{Handbook of Logic in Computer Science}, Vol. 4. Oxford University Press, 1995.

\end{thebibliography}

\end{document}


## IG's comments:

The introduction/contextualization, even for pure mathematicians, should be about the meaning of the theorem in the general framework (see current_track.md:interpretation)

For the extended abstract, I believe there are many terms that are thrown to the reader without any explanation what they are, or are the "homecooking" terms we've been using in our work but
  should be replaced by the term the mathematical community is expecting for the concept:
  "linear": has many possible meanings, is it linear as in vectors or linear logic, and is nonlinear curved then?)
  "point-generated actor-causal": there is not a single mention of actor-causal in the literature, this is a "homecooking" term we've been using in our work, it can't be thrown like that to a
  reader.
  "finite frame": as in Finite Frames: Theory and Applications edited by Peter G. Casazza, Gitta Kutyniok? If yes it probably should be referenced and contextualized, if not then it's worse, we're
  using a term to reference something else that its accepted
  "some event": very vague.
  "single-event point subframes": this could become understandable if the previous contextualization work has been done prior to it appearing in the text.
  "value category": absolutely not a recognized term in mathematics, it doesn't convey the meaning/points to the concept you want it to have, not witout prior explanation.
  "pointification": absolutely not a math term.
  "event-local values": could become understandable if 'event' and 'values' have been either a) replaced by the actual recognized math term they're supposed to refer to, or b) contextualized.
  "a value-level Benton linear--nonlinear adjunction": first occurence of Benton, again should be contextualized.
  "reflector": Is it 'The term cartesian reflector (A4.3.1) is used in the Elephant to say idempotent cartesian monad. This is because an idempotent monad is called a reflector (A1.1.1), there.'?
  Even if it is, this is a very obscure reference, we can't expect the readers to get it.
  "the sheafification hypothesis": what is it?
 
  For every single one of those, either:
  - change it to something that is a recognized math term, as in, it is an entry in the nLab.
  - give it context so that the reader knows what you're talking about.

  Also, why does the title say bridge and not adjunction, when that is a recognized math term we can rely on.